% Source: https://artofproblemsolving.com/community/c3h2707228_algebra_problem
% Stable identity: aops-topic-2707228-q01
The condition is $2p^2-4p-3=0$, whose discriminant $16+24=40$ is positive, so $p=1\pm\dfrac{\sqrt{10}}{2}$ is real. Using $2p^2=4p+3$, the numerator is $-4p$ and the denominator is $4(4p+3)-8p-3=8p+9$, so the expression equals $\dfrac{-4p}{8p+9}$, which depends on $p$. For $p=\dfrac{2+\sqrt{10}}{2}$:
\[\dfrac{-4-2\sqrt{10}}{17+4\sqrt{10}}=\dfrac{(-4-2\sqrt{10})(17-4\sqrt{10})}{129}=\dfrac{4-6\sqrt{10}}{43},\]
and for the other root the value is $\dfrac{4+6\sqrt{10}}{43}$. So the possible values are $\boxed{\dfrac{4\pm6\sqrt{10}}{43}}$.
